% Fichier LaTeX enrichi avec images UPDF
% Exercice 3 - Amérique du Nord ^
% 18 points

\documentclass[12pt]{article}
\usepackage[utf8]{inputenc}
\usepackage[T1]{fontenc}
\usepackage[french]{babel}
\usepackage{graphicx}
\usepackage{amsmath}
\usepackage{amssymb}
\usepackage{geometry}
\geometry{a4paper, margin=2cm}

\begin{document}

\section*{Exercice 3 (18 points)}

% Images de l'exercice
\begin{center}
\includegraphics[width=0.8\textwidth]{images/dnb_2025_06_ameriquenord_3_Image_007.gif}
\includegraphics[width=0.8\textwidth]{images/dnb_2025_06_ameriquenord_3_Image_008.gif}
\includegraphics[width=0.8\textwidth]{images/dnb_2025_06_ameriquenord_3_Image_009.gif}
\end{center}

\begin{enumerate}
\item Montrer que, lorsque le nombre choisi est 4, le résultat obtenu avec le programme A est 5.

\item Montrer que, lorsque le nombre choisi est $- 2$, le résultat obtenu avec le programme A est 5.

\item Justifier que l'affirmation suivante est vraie :

\begin{center}\og Le programme A donne toujours le même résultat. \fg\end{center}

\item Lorsque le nombre choisi est 10, quel résultat obtient-on avec le programme B ?

\item Il existe exactement deux nombres pour lesquels les programmes A et B fournissent à chaque fois des résultats identiques.

Quels sont ces deux nombres?

\end{enumerate}

\end{document}
